\documentclass[10pt]{report}
\usepackage[utf8x]{inputenc}
\usepackage{amsthm,amsmath,latexsym, amssymb}
\newtheorem{theorem}{Theorem}
\newtheorem*{theorem*}{Theorem}
\newtheorem*{proposition}{Proposition}
\newtheorem{definition}{Definition}
\newtheorem*{axiom}{Axiom}
\newtheorem*{notation}{Notation}
\newtheorem*{question}{Question}
\newtheorem*{lemma}{Lemma}
\newcommand{\cl}[1]{\overline{#1}}
%opening
\title{Topology}
\author{John W. Bales $<<$ Ben Fitzpatrick}
\date{}
\begin{document}

\maketitle

\begin{abstract} These are the definitions and theorems from Ben Fitzpatrick's
topology sequence MH 650, MH 651, MH 652 for the Auburn University academic year
1970-1971 as transcribed by his student John W. Bales. Much editorial assistance was provided by fellow student Sam
Creswell who made his notes available as well.

\end{abstract}

The following material is from Professor Ben Fitzpatrick's topology sequence taught at Auburn University in the 1970-71
academic year. The original numbering of the theorems has been maintained. Where Dr. Fitzpatrick inserted additional
theorems into the sequence we have retained the numbering which he gave to such extra theorems, if any. Although the
definitions were not originally numbered, we have numbered them in the order in which they were given in class. With
regard to Theorem Z, it was Dr. Fitzpatrick's custom to give it a designation outside the regular numbering sequence. 

\vspace{12pt}

``Point'' and ``region'' are undefined.\\

\begin{axiom}[0] Every region is a point set.\\
\end{axiom}

\begin{definition} The statement that the point $P$ is a limit point of the set
$M$ means that if $R$ is a region containing $P$ as an element, then $R$
contains a point of $M$ different from $P$.\\
\end{definition}

\begin{definition} The statement that the point $P$ is a boundary point of the
point set $M$ means that if $R$ is a region containing $P$, then $R$ contains a
point of $M$ and a point not in $M$.\\
\end{definition}

\begin{notation} If $M$ is a point set, then ${\cl{M}}$ denotes the point
set to which $P$ belongs if and only if $P$ is a point of $M$ or a limit point
of $M$.\\
\end{notation}

\begin{definition}: The statement that the point set $M$ is closed means that if
$P$ is a limit point of $M$, then $P$ belongs to $M$.\\
\end{definition}

\begin{theorem} No point of a region is a boundary point of that region.\\
\end{theorem}

\begin{theorem} If $M$ is a point set, then ${\cl{M}}$ is closed.\\
\end{theorem}

\begin{definition}
 An order relation on a set $M$ is a collection $\Omega$ of ordered pairs of
elements of $M$ with the following properties:
\begin{enumerate}
 \item If $x$ and $y$ are two elements of $M$, then $(x,y)$ or $(y,x)$ is an
element of $\Omega.$
 \item If $(x,y)$ and $(y,z)$ are elements of $\Omega$, then $(x,z)$ is an element
of $\Omega.$
 \item If $x$ is in $M$, then $(x,x)$ is not in $\Omega.$
\end{enumerate}
If $\Omega$ is an order relation on $M$, then the element $x$ of $M$ is said to
precede the element $y$ of $M$ according to $\Omega$ if and only if $(x,y)$ is
in $\Omega.$
\end{definition}

\begin{definition}
 A well-ordering on $M$ is an order relation $\Omega$ on $M$ such that if $H$ is
a subset of $M$, then there is an ``first'' element of $H$, i.e. an element $x$
of $H$ such that if $y$ is any other element of $H$, then $x$ precedes $y$
according to $\Omega.$
\end{definition}

\begin{definition}
 The collection $G$ of point sets is said to be monotonic if and only if it is
true that if $g$ and $h$ are two members of $G$, then either $g$ is a subset of
$h$ or $h$ is a subset of $g$.
\end{definition}

\begin{definition}
 The pointset $M$ is perfectly compact if and only if it is true that if
$\mathcal{H}$ is a monotonic collection of subsets of $M$, then there is a point
$P$ which is either common to all the members of $\mathcal{H}$ or is a limit
point of every member of $\mathcal{H}$.
\end{definition}

\begin{theorem}
 If $G$ is a collection of regions covering the closed and perfectly compact
point set $M$ and $\Omega$ is a well-ordering on $G$, then there exists a member
$y$ of $G$ such that the collection of regions in $G$ which $g$ does not
precede according to $\Omega$ covers $M$.
\end{theorem}

\begin{theorem*}[$\mathbb{Z}$]
 If $G$ is a non-degenerate set, then there exists a well-ordering on $G$.
\end{theorem*}

\begin{theorem}
 If $G$ is a collection of regions covering the closed and perfectly compact
point set $M$, then some finite sub-collection of $G$ covers $M$.
\end{theorem}

\begin{axiom}[$1_3$]
 There exists a sequence $G_1, G_2, G_3, \cdots$ such that
\begin{enumerate}
 \item if $n$ is a positive integer, then $G_n$ is a collection, each member of
which is a region, and $G_n$ covers $S$.
\item If $n$ is a positive integer, $G_{n+1}$ is a sub-collection of $G_n$.
\item if $R$ is a region, $A$ is a point of $R$ and $B$ is a point of $R$, then
there exists a positive integer $n$ such that if $g$ is a region of the
collection $G_n$ and $g$ contains $A$, then ${\cl{g}}$ is a subset of $R$
and does not contain $B$ unless $B$ is $A$.
\end{enumerate}

\end{axiom}

\begin{theorem}
 If $P$ is a point, then there exists a sequence $R_1,R_2,R_3,\cdots$ such that
\begin{enumerate}
 \item[(i)] if $n$ is a positive integer, then $R_n$ is a region containing $P$
 \item[(ii)] if $n$ is a positive integer, then $\cl{R}_{n+1}$ is a subset of $R_n$.
 \item[(iii)] $P$ is the only point common to all the terms of the sequence $R_1, R_2, R_3, \cdots$
 \item[(iv)] if $R$ is a region containing $P$, then there is a positive integer $n$ such that $R_n$
is a subset of $R$.
\end{enumerate}
\end{theorem}

\begin{definition}
 The sequence $R_1,R_2,R_3,\cdots$ in the previous theorem is said to ``close down'' on the point
$P$.
\end{definition}

\begin{theorem}
 If $P$ is a point and $G$ is a finite collection of regions each containing $P$, then there is a
region containing $P$ which is a subset of each member of $G$.
\end{theorem}

\begin{theorem}
 If a point $P$ is a limit point f the sum of two point sets $H$ and $K$, then it is a limit point of
one of them.
\end{theorem}

\begin{theorem}
 No finite point set has a limit point.
\end{theorem}

\begin{theorem}
 If the point $P$ is a limit point of the point set $M$, then there exists a sequence
$P_1,P_2,P_3,\cdots$ of points of $M$ all distinct from $P$ and from one another such that $P$ is
the sequential limit point of the sequence $P_1,P_2,P_3,\cdots$.
\end{theorem}

\begin{definition}
 The statement that the point $P$ is the sequential limit point of the point sequence
$P_1,P_2,P_3,\cdots$ means that if $R$ is a region containing $P$, then there is a positive integer
$N$ such that if $k$ is a positive integer greater than or equal to $N$, then $P_k$ is a point of
$R$.
\end{definition}

\begin{definition}
 The statement that the point set $M$ is bi-compact means that if $G$ is a collection of regions
covering $M$, then a finite sub-collection of $G$ covers $M$.
\end{definition}

\begin{theorem}
 Every bi-compact point set is closed and perfectly compact.
\end{theorem}

\begin{theorem}
 No point sequence has two sequential limit points and if the point $P$ is the sequential limit
point of the point sequence $P_1,P_2,P_3,\cdots$, then $P$ is the sequential limit point of every
infinite subsequence of $P_1,P_2,P_3,\cdots$.
\end{theorem}

\begin{theorem}
 If $\alpha$ is a point sequence and the set of all terms of $\alpha$ is compact, then some
subsequence of $\alpha$ has a sequential limit point.
\end{theorem}

\begin{definition}
 The statement that the point set $M$ is compact means that if $K$ is an infinite subset of $M$, then
$K$ has a limit point.
\end{definition}

\begin{theorem}
 If the set of all terms of the infinite point sequence $\alpha$ is compact and there do not exist
two points each of which is a sequential limit point of some subsequence of $\alpha$, then $\alpha$
has a sequential limit point.
\end{theorem}

\begin{definition}
The statement that the point set $D$ is a domain means that if $P$ is a point $D$, then there is a
region which contains $P$ and is a subset of $D$.
\end{definition}

\begin{theorem}
 If $H$ and $K$ are two mutually exclusive closed point sets and one of them is perfectly compact, then there exist
mutually exclusive domains $D_H$ and $D_K$ containing $H$ and $K$ respectively.
\end{theorem}

\begin{theorem}
 If $M$ is a point set and $G$ is a collection of regions covering $\cl{M}$, then there exists a
subset $K$ of $M$ with no limit point such that 
\begin{enumerate}
 \item[(1)] no region of $G$ contains two points of $K$ and 
 \item[(2)] if $P$ is a point of $M$, then some region of $G$ contains $P$ and intersects $K$.
\end{enumerate}
\end{theorem}

\begin{theorem}
 If $G$ is a collection of domains covering the bi-compact point set $M$, then some finite sub-collection of $G$ covers
$M$.
\end{theorem}

\begin{definition}
 A point set $M$ is said to be separable if and only if there is a countable subset $K$ of $M$ such that every point of
$M$ is a point of $K$ or a limit point of $K$. The point set $M$ is said to be completely separable if and only if there
is a countable collection $G$ such that each element of $G$ is a domain and such that if $P$ is a point of $M$ and $R$
is a region containing $P$, then some domain of the collection $G$ contains $P$ and is a subset of $R$.
\end{definition}

\begin{theorem}
 If $M$ is a point set, then the following statements are equivalent:
\begin{enumerate}
 \item $M$ is completely separable.
 \item Every subset of $M$ is separable.
 \item If $K$ is an uncountable subset of $M$, then $K$ has a limit points.
 \item If $G$ is a collection of domains covering $M$, then some countable sub-collection of $G$ covers $M$.
\end{enumerate}
\end{theorem}

\begin{theorem}
 If $H$ and $K$ are two mutually exclusive closed point sets and $G$ is a monotonic collection of closed sets and each
element of $G$ intersects both $H$ and $K$ and no element of $G$ is the sum of two mutually exclusive closed point sets,
one intersecting $H$ and the other intersecting $K$, and some element of $G$ is perfectly compact, then the common part
of all the elements of $G$ is closed and perfectly compact and intersects both $H$ and $K$ and is not the sum of two
mutually exclusive closed point sets, one intersecting $H$ and the other intersecting $K$.
\end{theorem}

\begin{definition}
 The statement that the point sets $H$ and $K$ are mutually separated means that they do not intersect and neither of
them contains a limit point of the other.
\end{definition}

\begin{definition}
 The point set $M$ is connected if and only if $M$ is not the sum of two mutually separated point sets.
\end{definition}

\begin{theorem}
 If the connected point-set $M$ is the subset of the sum of two mutually separated point sets, then it is a subset of
one
of them.
\end{theorem}

\begin{notation}
 If $G$ is a collection of sets, then $G^*$ denotes the set to which $P$ belongs if and only if $P$ is an element of
some set in $G$.
\end{notation}

\begin{theorem}
 If $G$ is a collection of connected point sets and some element of $G$ intersects every other element of $G$, then
$G^*$
is connected.
\end{theorem}

\begin{definition}
 If each of $M$ and $N$ is a point set, then the statement that $T$ is a transformation from $M$ onto $N$ means that $T$
is a collection each element of which is an ordered pair whose first term is an element of $M$ and whose second term is
an element of $N$ and no two elements of $T$ have the same first term and every element of $M$ is the first term of some
member of $T$ and every element of $N$ is the second term of some member of $T$.

The statement that the transformation from $M$ onto $N$ is continuous means that if $x$ is a point of $M$ and $\epsilon$
is a region containing $T(x)$, then there exists a region $\delta$ containing $x$ such that if $q$ belongs to the common
part of $\delta$ and $M$, then $T(q)$ is in $\epsilon$.
\end{definition}

\begin{theorem}
 If $T$ is a continuous transformation from $M$ onto $N$, then
\begin{enumerate}
 \item[(1)] $N$ is connected if $M$ is,
 \item[(2)] $N$ is bi-compact if $M$ is, and
 \item[(3)] $N$ is separable if $M$ is.
\end{enumerate}
\end{theorem}

\begin{theorem}
Every perfectly compact point set is compact and every compact point set is perfectly compact.
\end{theorem}

\begin{theorem}
 Every point of a non-degenerate connected point set is a limit point of that point set.
\end{theorem}

\begin{theorem}
 If $M$ is connected and $M$ is a subset of $L$ and $L$ is a subset of $\cl{M}$, then $L$ is connected.
\end{theorem}

\begin{theorem}
 If $H$ and $K$ are two point sets and $P$ is a boundary point of their common part, then $P$ is a boundary point of one
on them.
\end{theorem}

\begin{definition}
 If $M$ is a point set, then the statement that $C$ is a component of $M$ means that $C$ is a connected subset of $M$
and
is not a proper subset of any connected subset of $M$.
\end{definition}

\begin{theorem}
 If $M$ is a point set, then every point of $M$ belongs to a component of $M$.
\end{theorem}

\begin{theorem}
 If $M$ is a point set and $n$ is a positive integer such that $M$ has only $n$ components, then $M$ is the sum of $n$
mutually separated point sets
\end{theorem}

\begin{definition}
 A continuum is a closed and connected point set.
\end{definition}

\begin{theorem}
 If $G$ is a monotonic collection of perfectly compact continua, then the common part of the elements of $G$ is a
perfectly compact continuum.
\end{theorem}

\begin{definition}
 Suppose $H$ and $K$ are two mutually exclusive closed point sets and $M$ is a continuum which intersects both $H$ and
$K$, then $M$ is said to be irreducible from $H$ to $K$.
\end{definition}

\begin{theorem}
 If $M$ is a closed and perfectly compact point set that intersects both of the
two mutually exclusive closed point sets $H$ and $K$ and $M$ is not the sum of
two mutually exclusive closed point sets one intersecting $H$ and the other
intersecting $K$ but every closed proper subset of $M$ which intersects both
$H$ and $K$ is the sum of two such point sets, then $M$ is an irreducible
continuum from $H$ to $K$.
\end{theorem}

\begin{theorem}
 If the perfectly compact continuum $M$ intersects both of two mutually
exclusive closed point sets $H$ and $K$, then some sub-continuum of $M$ is
irreducible from $H$ to $K$.
\end{theorem}

\begin{definition}
 A topological transformation is a continuous and reversible transformation $T$
such that $T^{-1}$ is also continuous.
\end{definition}

\begin{definition}
 The point set $M$ is topologically equivalent to the point set $N$ if and only
if there is a topological transformation from $M$ onto $N$.
\end{definition}

\begin{theorem}
 Topological equivalence is an equivalence relation, i.e.
 \begin{enumerate}
  \item A point set is topologically equivalent to itself.
  \item If the point set $M$ is topologically equivalent to the point set $N$,
then $N$ is topologically equivalent to $M$.
  \item If the point set $M$ is topologically equivalent to the point set $N$
and $N$ is topologically equivalent to the point set $O$, then $M$ is
topologically equivalent to $O$.
 \end{enumerate}
\end{theorem}

\begin{theorem}
 Suppose $H$ and $K$ are two closed point sets and $f$ is a continuous
transformation from $H$ to a subset of the point set $M$ and $g$ is a
continuous transformation from $K$ to a subset of the point set $M$.
Furthermore, suppose that if $x$ belongs to $H$ and to $K$, then $f(x)=g(x)$.
Then $f\cup g$ is a continuous transformation whose initial set is $H\cup K$.
\end{theorem}

\begin{theorem}
 If $f$ is a continuous and reversible transformation whose initial set is
bi-compact, then $f^{-1}$ is also continuous.
\end{theorem}

\begin{theorem}
 Suppose $f$ is a transformation with initial set $M$ and $x$ is a point of $M$,
then $f$ is continuous at $x$ if and only if it is true that if
$P_1,P_2,P_3\cdots`$ is a point sequence which has $x$ as its sequential limit
point and each term of this sequence is in $M$, then the point sequence
$f(P_1),f(P_2),f(P_3),\cdots$ has $f(x)$ as its sequential limit point.
\end{theorem}

\begin{theorem}
 If $M$ is a connected point set and $T$ is a connected subset of $M$ and $M-T$ is the sum of two mutually separated
point sets $H$ and $K$, then $H+T$ is connected.
\end{theorem}

\begin{definition}
 The statement that the point or subset $P$ of the connected point set $M$ separates the point or subset $H$ of $M$ from
the point or subset $K$ of $M$ means that $M-P$ is the sum of two mutually separated point sets, one containing $H$ and
the other containing $K$.
\end{definition}

Note:
 In theorems \ref{T:S1}, \ref{T:S2} and \ref{T:S3}, $M$ denotes a connected point set and $P$, $A$, $B$ and $C$
represent points of $M$.

\begin{theorem}\label{T:S1}If $P$ separates $A$ from $B$ in $M$, then $P$ separates $B$ from $A$ in $M$ and $A$ does not
separate $P$ from $B$ in $M$.
\end{theorem}
\
\begin{theorem}\label{T:S2}If $P$ separates $A$ from $C$ in $M$ and $B$ separates $A$ from $C$ in $M$, then either $P$
separates $A$ from $\{B,C\}$ in $M$ or $B$ separates $A$ from $\{PmC\}$ in $M$.
\end{theorem}

\begin{theorem}\label{T:S3}If $P$ separates $A$ from $B$ in $M$ and $B$ separates $P$ from $C$ in $M$, then $B$
separates
$\{A,P\}$ from $C$ in $M$ and $P$ separates $A$ from $\{B,C\}$ in $M$.
\end{theorem}

[Note: In Dr. Fitzpatrick's 1970-1971 topology courses, he did not state a theorem 39 or 40.]\\

\setcounter{theorem}{40}

\begin{axiom}[1] Consists of Axiom $1_3$ plus
\begin{enumerate}
 \item[4.] if $M_1,M_2,M_3,\cdots$ is a simply infinite sequence, each term of
which is a closed point set, and each term of it contains the next term of it,
and for each positive number $N$, $M_N$ is a subset of the closure of some
element of $G_N$, then there is a point common to all the terms of the sequence
$M_1,M_2,M_3,\cdots$.
\end{enumerate}
\end{axiom}

\begin{theorem}
 If $M$ is a closed point set and every point of $M$ is a limit point of $N$, then $M$ is uncountable.
\end{theorem}

\begin{theorem}
 Every metric space satisfies Axiom $1_3$.
\end{theorem}

\begin{theorem*}[$42^\prime$]
 Every complete metric space satisfies Axiom 1.
\end{theorem*}

\begin{theorem}
 Every locally compact $1_3$ space satisfies Axiom 1.
\end{theorem}

\begin{definition}
 The statement that a space is locally compact means that if $P$ is a point of the space, then there is a compact
region containing $P$.
\end{definition}

\begin{theorem}
 No domain is the sum of countable many closed point sets, no one of which contains a domain.
\end{theorem}

\begin{theorem}
 No closed point set $M$ is the sum of countably many closed point sets such that if $T$ is any one of them every point
of $T$ is a limit point of $M_T$.
\end{theorem}

\begin{definition}
 Suppose $M$ is a set and $T$ is a set and $\Omega_1$ is a well-ordering on $M$ and $\Omega_2$ is a well-ordering on
$T$. The statement that $\Omega_1$ is ordinally similar to $\Omega_2$ means that there is a reversible transformation
$\phi$ from $M$ onto $T$ such that if $a$ and $b$ are two elements of $M$, then $a$ precedes $b$ according to $\Omega_1$
if and only if $\phi(a)$ precedes $\phi(b)$ according to $\Omega_2$.(i.e. there is a reversible order-preserving
transformation from $M$ onto $T$.)
\end{definition}

\begin{definition}
 The statement that the subset $M^\prime$ of $M$ is an initial (final) segment of $M$ according to $\Omega$ means that
there is an element $P$ of $M$ such that $M^\prime$ is the set of all elements of $M$ preceding $P$ (not preceding $P$)
according to $\Omega$.
\end{definition}

\begin{theorem*}[unnumbered] If $\Omega_1$ is a well-ordering on $M$ and $\Omega_2$ is a well-ordering on $T$, then
either (1) there is a reversible order-preserving transformation from $M$ onto $T$ or (2) there is a reversible
order-preserving transformation from $M$ onto an initial segment of $T$ or (3) there is a reversible order-preserving
transformation from an initial segment of $M$ onto $T$.
\end{theorem*}

\begin{definition}
 The statement that $A$ is cardinally equivalent to $B$ means that there is a reversible transformation from $A$ onto
$B$. The statement that $A$ is cardinally less than $B$ means that $A$ is cardinally equivalent to a subset of $B$ but
is not cardinally equivalent to $B$. The statement that $A$ is cardinally greater than $B$ means that $B$ is cardinally
less than $A$.
\end{definition}

\begin{definition}
 If $M$ is a connected point set and $P$ is a point of $M$, then $P$ is said to be a separating point of $M$ if $M-P$ is
not connected and is a non-separating point of $M$ is $M-P$ is connected.
\end{definition}

\begin{definition}
 An arc is a non-degenerate compact continuum with no more than two non-separating points.
\end{definition}

\begin{definition}
 The statement that $M$ is an arc from $A$ to $B$ means that $M$ is a compact continuum and if $P$ is a point of $M$
other than $A$ or $B$, then $P$ is a separating point of $M$.
\end{definition}

\begin{question}
 If $M$ is an arc are there two points $A$ and $B$ such that $M$ is an arc from $A$ to $B$?
\end{question}

\begin{question}
 Is there a non-degenerate compact continuum $M$ such that no point of $M$ is a non-separating point of $M$? That is,
is there a non-degenerate compact continuum $M$ such that every point of $M$ is a separating point of $M$?
\end{question}

\begin{proposition}[A1]
If $M$ is an arc than $M$ has two non-separating points.
\end{proposition}


\begin{proposition}[A2]
 If $M$ is an arc from $A$ to $B$ and $P$ is a  point of $M$ different from $A$ and from $B$, then $P$ separates $A$
from $B$ in $M$.
\end{proposition}

\begin{proposition}[A3]
 If $M$ is an arc from $A$ to $B$, $P$ is a point of $M$ different from $A$ and from $B$ and $H$ and $K$ are two
mutually separated point sets whose sum is $M-P$, then $A$ is in $H$ and $B$ is in $K$ or $A$ is in $K$ and $B$ is in
$H$.
\end{proposition}

\begin{proposition}[A4]
 Under the hypothesis, $H+P$ is an arc from $P$ to $A$ or to $B$.
\end{proposition}

\begin{proposition}[A5]
 If $M$ is an arc from $A$ to $B$, then $M$ is irreducible from $A$ to $B$.
\end{proposition}

\begin{proposition}[A6]
 If $M$ is an arc from $A$ to $B$ and $P_1,P_2,P_3,\cdots$ is an infinite sequence of points of $M$ and for each
positive integer $N$, $P_{N+1}$ separates the two-point set $\{A,P_N\}$ from $B$ in $M$, then $P_1,P_2,P_3,\cdots$ has a
sequential limit point which either is $B$ or separates $B$ from $\{A,P_1,P_2,P_3,\cdots\}$ in $M$.
\end{proposition}

\begin{theorem*}[$30^\prime$]
 If $M$ is a compact continuum and $K$ is a closed subset of $M$, then there is a sub-continuum of $M$ which contains
$K$ and is such that no proper sub-continuum of it contains $K$.
\end{theorem*}

\begin{theorem}
 If $M$ is a closed and perfectly compact point set and $H$ and $K$ are two mutually exclusive closed subsets of $M$
and $M$ contains no continuum which intersects both $H$ and $K$, then $M$ is the sum of two mutually exclusive closed
point sets, one containing $H$ and the other containing $K$.
\end{theorem}

\begin{theorem}
 If $H$ and $K$ are two mutually exclusive and closed point sets and $M$ is a compact continuum which is irreducible
from $H$ to $K$, then $M-M\cdot H$ is connected and $M-M\cdot(H+K)$ is connected and every point of $M\cdot H$ is a
limit point of $M-M\cdot H$.
\end{theorem}

\begin{theorem*}[unnumbered]
 If $H$ is a collection of regions covering $S$ and $M$ is a point set, $Z$ is a subset of $M$ and $G$ is a collection
of regions and for each point $P$ of $M$ there exists a region of the collection $G$ that contains every region of the
collection $H$ that contains $P$ and if every subset of $M$ that is cardinally greater than $Z$ has a limit point, then
some subset of $G$ cardinally less than or equivalent to $Z$ covers $M$.
\end{theorem*}

\begin{theorem}
 If $M$ is a closed and compact point set and $K$ is a component of $M$ and $H$ is a closed subset of $M-K$, then $M$ is
the sum of two mutually separated point sets one containing $K$ and the other containing $H$.
\end{theorem}

\begin{theorem}
 If $M$ is a continuum and $D$ is a domain with respect to $M$ and $D$ is a proper subset of $M$ and $\cl{D}$ is
compact and $K$ is a component of $D$, then the boundary of $D$ with respect to $M$ contains a limit point of $K$.
\end{theorem}

\begin{definition}
 If $M$ is a point set and $D$ is a subset of $M$, then the statement that $D$ is a domain with respect to $M$ means
that no point of $D$ is a limit point of $M-D$.
\end{definition}

\begin{definition}
 If $T$ is a subset of the point set $M$, then the statement that the point $X$ of $M$ is a boundary point of $T$ with
respect to $M$ means that every region containing $X$ contains a point of $T$ and a point of $M-T$.
\end{definition}

\begin{proposition}[A7]
 If $T$ is a continuous and reversible transformation whose initial set $M$ is an arc from $A$ to $B$, then $T$ is a
topological transformation whose final set $N$ is an arc from $T(A)$ to $T(B)$.
\end{proposition}

\begin{proposition}[A8]
 If $M$ is an arc from $A$ to $B$ and $P$ is a point of $M$ different from $A$ and from $B$ and $M-P$ is the sum of two
mutually separated point sets $H$ and $K$ and $H^\prime$ and $K^\prime$ are two mutually separated point sets whose sum
is $M-P$, then either $H^\prime$ is $H$ and $K^\prime$ is $K$ or $H^\prime$ is $K$ and $K^\prime$ is $H$.
\end{proposition}

\begin{proposition}[A9]
 Suppose $M$ is an arc from $A$ to $B$; $\{P_1,P_2,P_3,\cdots\}$ is a countable dense subset of $M-\{A,B\}$; $n_1=1$;
for each positive integer $j$, $n_{j+1}$ is the smallest integer $m>n_j$ such that one component of $M-P_m$ contains
$\{P_1,P_2,P_3,\cdots\,P_{n_j}\}$; and $P$ is a limit point of $\{P_1,P_2,P_3,\cdots\}$. Then $P$ is a non-separating
point of $M$.
\end{proposition}

\begin{definition}
 The statement that the point set $M$ is the limiting set of the sequence of point sets $M_1,M_2,M_3,\cdots$ means that
 \begin{enumerate}
  \item there exist a point $X$ such that if $R$ is a region containing $X$ and $i$ is a positive integer, then there is
a positive integer $j$ greater than $i$ such that $R$ intersects $M_j$ and
  \item $M$ is the set of all such points.
 \end{enumerate}
\end{definition}

\begin{theorem}
 If $M_1,M_2,M_3,\cdots$ is a sequence of connected point sets and $\cl{M_1+M_2+M_3+\cdots}$ is compact and for each
positive integer $N$, $a_N$ is a point of $M_N$ and the point sequence $a_1,a_2,a_3,\cdots$ has a sequential limit
point, then the limiting set of $M_1,M_2,M_3,\cdots$ is a compact continuum.
\end{theorem}

\begin{theorem}
 If $K$ is a closed proper subset of the continuum $M$, $\cl{M-K}$ is compact and $H$ is a component of $M-K$, then $K$
contains a limit point of $H$.
\end{theorem}

\begin{theorem}
 If $K$ is a component of a closed and compact point set $M$ and $T$ is a closed point set containing no point of $K$,
then 
\begin{enumerate}
 \item there is a domain $D$ containing $K$ such that no point of $T+M$ is a boundary point of $D$ and $D$ contains no
point of $T$ and
 \item $K$ is a component of $M+T$.
\end{enumerate}
\end{theorem}

\begin{definition}
 Suppose $M_1,M_2,M_3,\cdots$ is a simply infinite sequence each term of which is a point set. Then the statement that
the point set $M$ is the sequential limiting set of the sequence $M_1,M_2,M_3,\cdots$ means that $M$ is the limiting
set of every infinite sub-sequence of $M_1,M_2,M_3,\cdots$.
\end{definition}

\begin{theorem}
 If the sequence of point sets $M_1,M_2,M_3,\cdots$ has a completely separable limiting set, then some infinite
sub-sequence of that sequence has a sequential limiting set.
\end{theorem}

\begin{theorem}
 Every non-degenerate compact continuum contains two non-separating points.
\end{theorem}


\begin{theorem}
 If $M$ is a compact continuum and $H$ is a connected proper subset of $M$, then $M-H$ contains a non-separating point
of $M$.
\end{theorem}

\begin{theorem}
 If $M$ is an arc from $A$ to $B$, then $M-A$ and $M-B$ are connected.
\end{theorem}

\begin{theorem}
 If $M$ is an arc from $A$ to $B$ and $P$ is a point of $M$ different from $A$ and from $B$, then $M-P$ is the sum of
two mutually separated point sets and if $H$ and $K$ are two mutually separated point sets whose sum is $M-P$, then one
of $H$ and $K$ contains $A$ and the other contains $B$.
\end{theorem}

\begin{theorem}
 If $M$ is an arc from $A$ to $B$ and $P$ is a point of $M$ different from $A$ and from $B$ and $M-P$ is the sum of two
mutually separated point sets $H$ and $K$ and $A$ is in $H$, then $H+P$ is an arc from $A$ to $P$ and $K+P$ is an arc
from $B$ to $P$.
\end{theorem}

\begin{definition}
 The statement that $P$ precedes $Q$ in the order from $A$ to $B$ means that either $P$ is $A$ and $Q$ is $B$ or $P$
separates $A$ from $\{Q,B\}$.
\end{definition}

\begin{definition}
 The statement that $T$ is a segment from the point $P$ of the arc $M$ to the point $Q$ of the arc $M$ means that $T$
is the set of points to which the point $C$ belongs only in case $C$ is in $M$ and $C$ precedes $Q$ in $M$ and $P$
precedes $C$ in $M$ (both in the order from $A$ to $B$ in $M$ or both in the order from $B$ to $A$ in $M$ provided $M$
is an arc from $A$ to $B$.)
\end{definition}

\begin{definition}
 If $M$ is an arc from $A$ to $B$ the statement that $T$ is a sect containing $A$ means that there is a point $P$ of
$M$ different from $A$ such that $T$ is the subset of $M$ to which the point $Q$ belongs only in case $Q$ precedes $P$
in the order of $A$ to $B$ in $M$.
\end{definition}

\begin{theorem*}[58.5]
 If $M$ is an arc and $T$ is a subset of $M$, then $P$ is a limit point of $T$ only in case either
 \begin{enumerate}
  \item $P$ is a separating point of $M$ and every segment of $M$ containing $P$ contains a point of $T$ different from
$P$ or
  \item $P$ is a non-separating point of $M$ and every sect of $M$ containing $P$ contains a point of $T$ different
from $P$.
 \end{enumerate}
\end{theorem*}

\begin{lemma}
 If $M$ is an arc from $A$ to $B$ and $C$ is a countable dense subset of $M$ which contains neither $A$ nor $B$ and $Z$
is a countable dense subset of the number segment $(0,1)$, then there is a reversible transformation $T$ from $C$ onto
$Z$ which preserves order.
\end{lemma}

\begin{theorem}
 Every arc is topologically equivalent to the number interval $[0,1]$.
\end{theorem}

\begin{definition}
 A simple closed curve is a non-degenerate compact continuum each two point subset of which separates it.
\end{definition}

\begin{theorem}
 No point of a simple closed curve is a separating point of it.
\end{theorem}

\begin{theorem}
 If $A$ and $B$ are points of the simple closed curve $J$, then $J$ is the sum of two arcs from $A$ to $B$ which have in
common only $A$ and $B$.
\end{theorem}

\begin{theorem}
 If $M$ is a connected point set, then there is no domain $D$ such that the common part of $D$ and $M$ is separable and
contains a point set $H$ which contains an uncountable point set $K$ such that if $X$ is any point of $K$, then $M-X$ is
the sum of two mutually separated point sets, one of which contains $H-X$.
\end{theorem}

\begin{definition}
 The statement that the point set $M$ is connected im kleinen at the point $P$ of $M$ means that if $U$ is an open set
containing $P$, then there is an open set $V$ containing $P$ such that if $qQ$ is a point of the common part of $V$ and
$M$, then some connected set of the common part of $U$ and $M$ contains both $P$ and $Q$.
\end{definition}

\begin{definition}
 The statement that the point set $M$ is open means that $M$ is a domain.
\end{definition}

\begin{theorem}
 If $M$ is connected im kleinen and $D$ is a domain with respect to $M$, then every component of $D$ is a domain with
respect to $M$.
\end{theorem}

\begin{definition}
 The statement that the collection $H$ of point sets is coherent means that if $H_1$ and $H_2$ are two sub-collections
of $H$ whose sum is $H$, then some member of $H_1$ intersects some member of $H_2$.
\end{definition}

\begin{theorem}
 If $G_1,G_2,G_3,\cdots$ is an Axiom 1 sequence for $S$ and for each positive integer $n$, $H_n$ is a finite coherent
sub-collection of $G_n$ and for each $n$ the closure of each element of $H_{n+1}$ is contained in some element of
$H_n$, then the common part of $H_1^*,H_2^*,H_3^*,\cdots$ is a compact continuum.
\end{theorem}

\begin{theorem}
 If $S$ is connected and connected im kleinen and $A$ and $B$ arer two points of $S$, then $S$ contains an arc from $A$
to $B$.
\end{theorem}

\begin{theorem*}[$64^\prime$]
 If $M$ is a connected point set and $G$ is a domain collection covering $M$ and $A$ and $B$ are two points of $M$, then
there is a finite coherent sub-collection of $G$ such that $G^*$ contains both $A$ and $B$.
\end{theorem*}

\begin{theorem}
 If the continuum $M$ is locally compact at the point $P$ of $M$ but is not connected im kleinen at $P$, then there is a
region $R$ containing $P$ and a sequence $K_1,K_2,K_3,\cdots$ of distinct components of $M\cdot\cl{R}$ such that each
$K_n$ contains a point of $R$ and a point on the boundary of $R$ and $K_1,K_2,K_3,\cdots$ converges to a continuum $K$
which contains $P$ and intersects the boundary of $R$.
\end{theorem}



\end{document}
